\begin{table}[htbp]
\centering
\caption{Standardized Effect Sizes}
\label{tab:sde}
\begin{tabular}{lcccccc}
\hline\hline
Outcome & $\hat{\beta}$ & SE & SD($Y$) & SDE & SE(SDE) & Classification \\
\hline
\multicolumn{7}{l}{\textit{Panel A: Pooled}} \\
Log Revenue & -0.0396 & 0.0635 & 0.9446 & -0.0419 & 0.0672 & Small negative \\
Rev.\ Growth & -0.0068 & 0.0483 & 0.4573 & -0.0149 & 0.1057 & Small negative \\
Log Expenses & -0.0426 & 0.0729 & 1.1510 & -0.0370 & 0.0633 & Small negative \\
\hline
\multicolumn{7}{l}{\textit{Panel B: Heterogeneous (by sector)}} \\
Arts/\allowbreak Culture/\allowbreak Education & 0.1142 & 0.0987 & 0.4830 & 0.2364 & 0.2044 & Large positive \\
Health/\allowbreak Human Services & -0.8460 & 0.6327 & 1.8667 & -0.4532 & 0.3389 & Large negative \\

\hline\hline
\end{tabular}
\begin{tablenotes}
\item \textit{Notes:} \textbf{Country:} United States. \textbf{Research question:} Does the IRS Form 990-EZ gross receipts threshold at \$200K constrain revenue growth for nonprofits near the compliance boundary, compared to similar-sized organizations between the old and new thresholds? \textbf{Policy mechanism:} The 2010 reform raised the Form 990-EZ eligibility ceiling from \$100K to \$200K, creating a new compliance discontinuity where organizations exceeding \$200K in gross receipts must file the full twelve-page Form 990 rather than the four-page 990-EZ, imposing incremental reporting costs on governance, compensation, and program activities. \textbf{Outcome definition:} Log annual gross receipts from IRS Form 990/990-EZ electronic filings, measuring total organizational revenue. \textbf{Treatment:} Binary; organization classified as near the \$200K threshold (mean gross receipts \$170K--\$220K in 2011--2013) versus mid-range control (\$120K--\$160K). \textbf{Data:} IRS Exempt Organizations Business Master File and ProPublica Nonprofit Explorer API, 2011--2022, organization-year panel, 508 organizations. \textbf{Method:} Two-way fixed effects DiD with organization and year fixed effects; standard errors clustered at the organization level. \textbf{Sample:} 501(c)(3) organizations with at least two baseline-period filing years (2011--2013) and mean gross receipts between \$120K and \$220K; excludes organizations with gross receipts above \$2M. SDE $= \hat{\beta} / \text{SD}(Y)$ where SD($Y$) is the pre-treatment standard deviation. Classification refers to magnitude, not statistical significance: Large ($|$SDE$| > 0.15$), Moderate ($0.05$--$0.15$), Small ($0.005$--$0.05$), Null ($< 0.005$).
\end{tablenotes}
\end{table}
